\subsection*{从保证定位权威来源}

一项保证先指定被计数的对象，再说明它由哪一项决定产生。表~\ref{tab:ch15-guarantee-path}~把原身份、结果和恢复责任接到各自的保存位置。购买库能否回退、平台保护多久、消息怎样重投，分别影响相应行的条件。

\begin{table}[htbp]
\centering\small
\caption{核心保证、权威来源与验证入口}\label{tab:ch15-guarantee-path}
\begin{tabularx}{\linewidth}{@{}p{28mm}X X@{}}
\toprule
保证对象 & 保存来源与允许迁移 & 推导和检验的位置\\
\midrule
同 S 原购买 & 原输入、P 与结果共同提交；复制和快照携带同一应用边界 & 第三至六、十一章，响应丢失、异参及切换历史\\
结果回收后的旧身份 & 活动会话保留准确身份依据，关闭后由持久资格拒绝旧入口 & 第十二章，空洞、迟到及缺失会话历史\\
原 C 的付款 & 固定范围和平台键，C/T 保留未知责任，可信证据形成唯一结果 & 第七、八、十二章，成功丢响应与下游窗口历史\\
旧库恢复后的原 O/EV & J 保存 $A_J/R_J$，物化事务比较完整目标并推进覆盖边界 & 第十四章及本章恢复演练，损坏行与零行回滚\\
原 B 的初次履约 & Q 共同保存 I/B、实际 F、原接收结果与 dispatch 待办 & 第九、十章，发布重试、消费提交后失联\\
多项部分退款 & J 保存原 RF、来源、金额、成功与未决占用，未知不释放额度 & 本章退款协议、归纳、有限模型与并发预占\\
\bottomrule
\end{tabularx}
\end{table}

本地 v1 的原成交余额与当前余额继续分别保存。例如 S1 当时为 7000，S2 之后当前为 4000，恢复 S1 返回前者。外部 v2 从独立平台账本核对资金，购买库的成功结果则固定 P/C/O/X 和原金额。两个动作引用不同资金对象，升级读取器按原版本恢复各自的结果。

\subsection*{保存迁移和执行迁移各自怎样继续}

复制接班延续已提交前缀，快照把原输入、绑定、会话与结果装在同一边界。旧备份恢复则先确定哪份权威历史继续存在，再按相同编号重建投影。物化事务遇到异归属、损坏结果或事件就整批回滚，本批游标保持，核查责任在事务外保存。提交未知以后重读原对象和边界，恢复者据实际提交事实继续。

跨域执行依靠另一条接力关系：任务先持久保存原身份，发送者按下游当前允许的方式查询或发送，可信结果再形成固定决定。消费者提交 F501 后 ACK 丢失，重投从 I/B 返回原 F501。外部 dispatch 已发送却未获回执，发送者沿原 dispatch 查询或按受保护规则重发。普通调用无法给出终态时，核查责任继续保存在对应任务中。

退款同时使用这两条关系。来源 C1/X1 给出可退收款，新的 RF 表达新的退款意图，未知占用承担可能已经成功的金额。成功把占用转为已退，足够最终拒绝才释放占用。购买库回退以后，退款投影从权威 RF 和原结果恢复，额度判断使用完整权威记录。

\subsection*{验证材料怎样与本章模型对应}

证明沿每种允许迁移检查关系，有限模型枚举指定身份、金额和阶段，数据库历史检查行锁、唯一性、完整比较及异常回滚。受控平台按明确的身份与查询规则产生资金事实，复制组件按自己的日志、提交和快照协议延续状态。各类证据分别指向它实际检验的模型。

本章的关系组合采用单逻辑购买分区，J、购买投影、平台与消费者分别提交。目录和 J 共享权威仲裁，首次受理先取得其开放资格。第十一章的主模型直接按完整 S 保存 M，配套复制实验采用协议客户端与序号的扩展编码，并随快照保存 U。实验中的结果投影与会话表示是对应接口的简化，完整字段的恢复责任按正文的 $R(S)$、会话状态和原事件定义检验。字段映射、实际运行条件、原始历史与复现入口集中在配套材料，读者据此判断每份证据覆盖哪项迁移。

集合合并承担已形成信息的传播。标签并集吸收重复标签，OR-set 用来源标记和因果上下文区分删除与并发新增，履约查询从已核验事实计算阶段。首次购买的余额判断、退款的来源额度判断和首次发货的授权判断，在各自稳定身份处形成决定，随后再传播对应事实。

接口最终对应 Alice 的具体购买：E1/K1 取得原 P1/C1，原平台范围和 \code{pay-C1} 指向 X1，原订单 O1 与事件 EV1 连接实际 F501。两项合法退款各取得自己的 RF 和原结果。E1 关闭后，旧提交入口结束准入，服务端任务与授权查询继续沿这些关联核对已经接受的工作。
\FloatBarrier
